\section{An independent executor: the tape machine and the spinor tape}\label{sec:tape}

The interval recipe of Section~\ref{sec:requests} has exactly one implementation, so it can only be checked against a hand calculation.
For one instruction, QAM goes further: the same plan is executed by several implementations, including a hardware core with two controllers, and for the hardware the traces are compared step by step. The instruction is a tape machine, the simplest model of a computer with a
separately stated plan. This section first describes the machine and the comparison, and then a hypothesis about where such a tape
might live, which shows how QAM records a claim whose conditions are far from discharged.

\subsection{Plan and execution}

The instruction \texttt{RUN\_TAPE} takes two objects from the request context. The \emph{program} is an operator object with a set of
control states, a halting state and a deterministic, possibly partial transition table (at most 64 states and 128 rules): for a control
state and a bit read from the tape, the bit to
write, the head movement ($-1$, $0$ or $+1$) and the next control state. The \emph{initial state} is a state object with the tape, the
head position and the control state. The request sets the tape mode, either fixed or growing (a blank cell is appended at the right edge
when needed), bounds of at most 1024 steps and 1024 cells, and optionally a full trace.

The result separates two questions that are easily confused:
\begin{itemize}[leftmargin=1.5em]
\item \emph{Was the request allowed to run?} This is the executor's outcome, \cont, \guard\ or \reject, as in Section~\ref{sec:requests}.
\item \emph{How did the computation end?} This is the machine's stop reason: \textsf{halted} (the halting state was reached),
\textsf{stuck} (no rule for the current state and bit), \textsf{boundary} (a step would leave the tape at its left end or, in fixed mode, at its right end; the whole step is discarded),
\textsf{step\_limit} or \textsf{cell\_limit} (the budget is exhausted).
\end{itemize}
An exhausted budget is explicitly not a verdict of non-termination. The machine does not claim that the computation would never end; it
reports that it did not end within the budget and returns the tape, head and state from which a new run with a larger budget can
continue. This is Turing's boundary in practical form: whether a program halts cannot be decided in general, so the executor never
pretends to decide it.

\subsection{Several executors for one plan}

The plan is executed in full by the Python adapter (the reference) and by a core written in SystemVerilog with two selectable
controllers. A Wolfram Language script and a Qiskit version, which builds every step as a reversible circuit over all decision cases,
implement the transition rule without the step and cell budgets; the Qiskit version has only the fixed tape and at most 18 qubits. The first controller addresses the
matching table entry directly. The second reads the whole table cyclically on every machine step, remembers the matching entry and applies
it after the last slot. Both share the logic for writing, head movement, tape growth and stop signals; the example instance has eight physical cells and six
table entries (three control states, two bits read). The Wolfram and Qiskit versions agree with the reference on the example traces that were run (two for Wolfram, one for Qiskit) and
with each other on all 24 start configurations of a two-cell program; the systematic comparison concerns the hardware.

A comparison suite runs both controllers in a hardware simulator and compares them with the reference traces obtained through the
ordinary executor. The suite covers 209 cases, among them 204 start configurations formed by all binary tapes of length one to three,
all head positions, three control states and both tape modes of a branching program. Both controllers agreed with the reference in all
209 cases, in stop reason and in every visible machine step. Clock counts differ as expected: the example run that adds one to the
binary number seven halts after four machine steps, which take four clock cycles on the direct controller and 24 on the cyclic one.
The suite is not only shown to pass. In a temporary copy of the hardware core, the tape was deliberately written also on a
\textsf{boundary} stop; the comparison then reported the first deviating case and step, showing that it detects this kind of error.

What the comparison establishes is limited and precise. It checks the executable transition semantics of the plan on a finite set of
cases. It does not interpret the machine, and the binary alphabet and the growth rule are model choices.

\subsection{The spinor tape hypothesis}

Why a tape machine, rather than any other test program? The research project in which QAM was developed pursues a question that this
article does not try to answer: whether a particular mathematical structure, built from spinors and their transport, carries a
computing machine of its own, a machine that runs not on hardware but in a mathematical carrier. We describe the hypothesis because it is
a good example of a claim that QAM must keep conditional for a long time, and because \texttt{RUN\_TAPE} is its executable counterpart.

In that structure there is a distinguished line along which spinors (two-component objects on which rotations act through $SU(2)$)
can be transported while satisfying the structure's admissibility conditions. The \emph{spinor tape hypothesis} states that this
line can serve as an extendable read/write support: finite stretches of a tape can be prolonged along
it without violating admissibility. Assuming this, the project derives that such stretches can be written and read and that a transport
step acts as a head movement. The motivation is structural. Transport along a line is
sequential and one-dimensional, and in Turing's model the tape is the simplest store with local access.

The accompanying \emph{conditional Z3 hypothesis} is named after Zuse's Z3. The Z3 had no conditional branch instruction; Rojas showed
that it is nevertheless universal in principle, by emulating branching and indirect addressing arithmetically, as universal as a computer
with bounded address space and at the price of very long programs~\cite{rojas98}. The hypothesis is named in this spirit (the link to Rojas's analysis is our gloss, not part of its statement): once branching,
persistent storage and return are available on the spinor carrier, the request-and-executor stack described here can emulate an
ordinary sequential control core.

The project states a local operational semantics for such a machine: a state $\Sigma=(\tau,h,q,\rho,\chi)$ consisting of a finite tape
section, a head position, a control state, a readout and a certificate of admissibility, together with primitive instructions that move
the head, produce branch predicates, update the control state or halt. The route assumes, besides the two hypotheses, that the structure's components
really play the computational roles assigned to them. Universality is \emph{not} proved, and the project says so. It isolates a
conditional route with five obligations, listed here in the recommended order in which each reduction assumes the previous ones:
\begin{enumerate}[label=\arabic*.,nosep,leftmargin=2em]
\item the control layer composes stably across neighbouring windows of the carrier;
\item the branch predicates are executable and remain decidable under composition;
\item return and looping can be resumed under repeated execution;
\item the tape is globally, not just locally, prolongable;
\item a standard sequential reference machine is explicitly simulated.
\end{enumerate}
Each has been reduced to a stated certificate (on finite windows for the first three, along chains of windows for the last two), but
only under the assumption that such a certificate exists and remains stable when windows are joined or chains are extended. None of the five is discharged.

\subsection{How QAM records it}

In a QAM record, ``the carrier is universal'' would therefore be an entry with status \st{conditional} whose open conditions are these five
obligations. By Theorem~\ref{thm:nopromotion}, nothing derived from it can be stronger, and by the impact query of
Section~\ref{sec:graph}, every result that uses it can be found when one of the obligations is settled.

The successful runs of \texttt{RUN\_TAPE} do not discharge any of the obligations. They check something else: that the declared plan is
executed correctly. A deterministic binary machine with a one-sided unbounded tape is universal in the idealised sense; with at most 1024
steps and cells per request the executor is, like the Z3 in Rojas's analysis, universal only in the bounded sense. That the tape can be
identified with the spinor carrier remains a premise that no run fulfils. If the carrier were universal, the halting problem for its
programs would be undecidable. That step and cell limits are recorded as stop reasons of their own and never read as ``does not halt'' is
therefore not a convenience. A single bounded request is a finite computation, but a run can be continued from its returned state
again and again, and for such continued runs no executor could decide halting. This is the very boundary the hypothesis aims to reach.
